\chapter{Conclusions and Future Work}

\section{Main Findings}
\label{sec:main_findings}

The results exposed in Chapter~\ref{chap:experimental_results} provide a clear answer to the research question formulated in Section~\ref{sec:research_gap}:\textbf{\textit{ modifying the glycemic ranges distribution of the training data can substantially improve glucose prediction in underrepresented clinically critical regions}}. This conclusion was reproduced in DiaTrend, REPLACE-BG, and T1DiabetesGranada, despite the particular differences between them, explored in Chapter~\ref{chap:eda}.

\begin{itemize}
    \item  The exploratory data analysis demonstrated why this intervention was necessary. Severe hypoglycemia represented less than \(1\%\) of the available glucose measurements in all three datasets, while normoglycemia appeared between 60 and 110 times -- depending on the dataset -- more frequently. Also, it showed that this imbalance was not caused only by a small group of highly represented patients. It was also observed at the patient level, confirming that the limited representation of hypoglycemia is an intrinsic characteristic of real-world CGM data. This finding directly illustrated the limitation of evaluating glucose prediction models simply through a global regression loss. Since most training windows belong to TIR or moderate hyperglycemia, the model "plays it safe" by assuming all the measurements are in the normoglycemic range, in order to reduce its average error. From a purely mathematical perspective, this behaviour may appear reasonable. But from a clinical perspective, it means that a model can obtain satisfactory global performance while almost systematically failing in the regions where an accurate prediction is much more critical.
    \item By adjusting the glycemic ranges distribution to achieve a fully symmetrical balance, this behaviour changed enormously. The primary statistical outcomes, \(A+B_{\mathrm{TBR_2}}\) and \(A+B_{\mathrm{TBR_1}}\), increased substantially across all three datasets. For TBR$_2$, the mean improvements were \(63.75\), \(76.81\), and \(75.93\) percentage points in DiaTrend, REPLACE-BG, and T1DiabetesGranada, respectively. For TBR$_1$, the corresponding improvements were \(66.58\), \(78.14\), and \(70.68\) percentage points. The one-sided paired Student's \(t\)-tests rejected the null hypothesis for both ranges in all three datasets, and these conclusions remained unchanged after the Holm correction.
    \item The numerical magnitude of these effects is as important as their statistical significance. ¡The aggregate TBR$_2$ A+B percentage increased from \(3.74\%\) to \(72.68\%\) in DiaTrend, from \(0.01\%\) to \(76.89\%\) in REPLACE-BG, and from \(13.24\%\) to \(88.87\%\) in T1DiabetesGranada. Similar improvements were obtained for TBR$_1$, where A+B reached \(70.35\%\), \(78.49\%\), and \(79.06\%\), respectively. These gains were mainly due to a large movement of predictions from Zone D into Zone A.
    \item Full balancing produced an important transformation in models behaviour. The unbalanced original models were rarely capable of producing clinically acceptable predictions in hypoglycemia, whereas the fully balanced models correctly represented a large majority of them, and across three independent datasets.
    \item Additionally, the intermediate semi-balanced configuration also provided valuable evidence. In every dataset, its results were located between those of normalized and full balancing, which suggest a progression pattern -- \textit{not proven in this study.}
        \begin{equation}
        AB_{\mathrm{TBR}}^{(\mathrm{normalized})}
        <
        AB_{\mathrm{TBR}}^{(\mathrm{semi})}
        <
        AB_{\mathrm{TBR}}^{(\mathrm{full})}.
        \end{equation}
    Nevertheless, semi-balancing approach was only investigated as an intermediate experimental condition rather than as the final preferred strategy.
    \item Another important finding is that the balancing target had a much greater influence than the oversampling method. Random resampling, SMOTER, and SMOGN produced very similar results for the same target distribution. In some cases, even simple random resampling produced larger improvements than other more complex synthetic methods. This result was not formally demonstrated.
    \item The improvements in hypoglycemia were accompanied by some numerical trade-offs. For example, full balancing increased global RMSE by \(3.23\) mg/dL in DiaTrend, \(4.70\) mg/dL in REPLACE-BG, and \(2.30\) mg/dL in T1DiabetesGranada. Similar deteriorations occured in TIR. However, these changes must be interpreted according to where the additional error occurred. TIR A+B remained at \(99.01\%\) in DiaTrend, \(99.46\%\) in REPLACE-BG, and \(99.39\%\) in T1DiabetesGranada. No TIR predictions entered Clarke Error Grid Zones D or E. Therefore, these deteriorations only reflect a reduction in numerical precision but within clinically acceptable regions, rather than the introduction of clinically dangerous errors.
    \item At the global level, clinical acceptability actually improved in all three datasets despite the increase in RMSE. Global A+B increased from \(95.23\%\) to \(96.09\%\) in DiaTrend, from \(95.08\%\) to \(97.55\%\) in REPLACE-BG, and from \(94.26\%\) to \(97.53\%\) in T1DiabetesGranada. At the same time, the global proportion of Zone D predictions decreased considerably. These results demonstrate that a higher global RMSE does not necessarily indicate a clinically worse model.
    \item The most important finding of this study, from a clinical perspective, is that glucose prediction is not an abstract regression problem in which every error has the same consequence. Behind each prediction is a person who may need act before losing the ability to respond independently. The practical difference between detecting and failing to detect a severe hypoglycemic event may actually be the difference between an early manageable intervention and a situation that leads to physical impairment, seizure, loss of consciousness, or even death. For this reason, even a small improvement in this context would be worthy of serious consideration. However, the improvements observed in the present experiments were far from modest: the proportion of clinically acceptable hypoglycemic predictions increased by approximately 64 to 78 percentage points across TBR$_2$ and TBR$_1$ in the primary fold-level analyses, and reached between \(70.35\%\) and \(88.87\%\) in the aggregate fully balanced results.

\end{itemize}

These results should not be interpreted as direct evidence that the evaluated models would actually perform satisfactorily in the real world. Such a conclusion would require prospective clinical validation under realistic deployment conditions, which is definitely not the aim of this study. However, the experiments clearly demonstrate that data balancing approaches can substantially reduce the type of predictive failures that would prevent a model from recognizing critical hypoglycemia events.

The broader conclusion is that glucose prediction models should not be optimized solely according to aggregate RMSE. When clinical impact is treated as a main priority, models can be substantially more useful and safer, even if they are slightly less accurate on average.  Consequently, future glucose prediction systems should evaluate both the magnitude and the clinical location of their errors, and should explicitly consider the training-data glycemic ranges distribution as a fundamental decision.

\section{Future Work}
\label{sec:future_work}

The results obtained in this thesis also open some possible directions for future research.

First, the target distribution could be extended beyond symmetrical full balancing. Equal representation of the five glycemic ranges was used in this study as a clear and naive approach, but it is not necessarily the optimal distribution. Future experiments could explore asymmetric targets that assign even greater presence to severe hypoglycemia. This would make it possible to identify the most clinically efficient distribution.

Second, the methodology should be evaluated with alternative predictive architectures: history lengths, prediction horizons, learning algorythms, etc.

Finally, data balancing could be combined with clinically aware loss functions that assign greater penalties to dangerous prediction errors. This would allow evaluating whether data-balancing and optimized loss functions are complementary to produce improvements in critical glycemic ranges.